% GENERATED FILE. Edit canonical lessons, then run npm run generate:books.
% canonical-source-hash: fnv1a64:39a67503ce87c26f
% canonical-lessons: JA-C509-takarakuji, JA-C509-kujibiki, JA-C509-chuusen-suru, JA-C509-kakuritsu, JA-C509-ateru, JA-R509-review-504-509

\chapter{Lotteries and Odds}
\label{ch:ja-lotteries-and-odds}
\hlchaptermodality{509}

\begin{chapteropening}
\textbf{By the end of this chapter:} \emph{I can talk about lotteries and chances.}

Complete the last lesson of chapter 509: I can talk about lotteries and chances.
\end{chapteropening}

\section[\texorpdfstring{\ja{たからくじ} (takarakuji)}{たからくじ (takarakuji)}]{\ja{たからくじ} (takarakuji) — a lottery; a lottery ticket}

\label{lesson:JA-C509-takarakuji}



\begin{quote}
Before the new one: say the Japanese for luckily, then the Japanese for fate.
\end{quote}

\subsection*{You'll want to know: \ja{たからくじ}}
\textbf{\ja{たからくじ}} — \emph{takarakuji} — \textquotedblleft{}a lottery; a lottery ticket\textquotedblright{}.

\begin{quote}
\textbf{\ja{もし} \ja{たからくじが} \ja{あたったら、うちを} \ja{かいます。}} — \emph{moshi takarakuji ga atattara, uchi o kaimasu} — \textquotedblleft{}If I win the lottery, I will buy a house.\textquotedblright{}
\end{quote}

\begin{cousinweb}
\textbf{\ja{たからくじ}} (\emph{takarakuji}) is a native Japanese compound (wago), inherited rather than borrowed from Chinese: built from \emph{takara}, \textquotedblleft{}treasure\textquotedblright{} and \emph{kuji}, \textquotedblleft{}lot\textquotedblright{} (JMdict, KANJIDIC2).
\end{cousinweb}

\begin{grammarlens}[title={What you've built}]
One more word for saying what happens if.
\end{grammarlens}

\subsection*{Your turn}
\begin{itemize}
\raggedright
  \item \emph{Say it:} \emph{takarakuji}
  \item \emph{Say it:} \emph{takarakuji}, once more
  \item \emph{Say it:} \emph{moshi takarakuji ga atattara, uchi o kaimasu} — If I win the lottery, I will buy a house.
  \item \emph{Recall:} say the Japanese for luckily, then the Japanese for fate, then say \emph{takarakuji} again
\end{itemize}

\begin{tcolorbox}[breakable,colback=teal!4,colframe=teal!35!black,title={Before you move on}]
What does \ja{たからくじ} mean? (\textquotedblleft{}A lottery; a lottery ticket\textquotedblright{}.) Say it once more.
\end{tcolorbox}
\section[\texorpdfstring{\ja{くじびき} (kujibiki)}{くじびき (kujibiki)}]{\ja{くじびき} (kujibiki) — drawing lots}

\label{lesson:JA-C509-kujibiki}



\begin{quote}
Before the new one: say the Japanese for fate, then the Japanese for a lottery.
\end{quote}

\subsection*{You'll want to know: \ja{くじびき}}
\textbf{\ja{くじびき}} — \emph{kujibiki} — \textquotedblleft{}drawing lots\textquotedblright{}.

\begin{quote}
\textbf{\ja{ばしょは} \ja{くじびきに} \ja{よります。}} — \emph{basho wa kujibiki ni yorimasu} — \textquotedblleft{}The places depend on drawing lots.\textquotedblright{}
\end{quote}

\begin{cousinweb}
\textbf{\ja{くじびき}} (\emph{kujibiki}) is a native Japanese compound (wago), inherited rather than borrowed from Chinese: built from \emph{kuji}, \textquotedblleft{}lot\textquotedblright{} and \emph{biki}, \textquotedblleft{}pulling\textquotedblright{} (JMdict, KANJIDIC2).
\end{cousinweb}

\begin{grammarlens}[title={What you've built}]
One more word for saying what happens if.
\end{grammarlens}

\subsection*{Your turn}
\begin{itemize}
\raggedright
  \item \emph{Say it:} \emph{kujibiki}
  \item \emph{Say it:} \emph{kujibiki}, once more
  \item \emph{Say it:} \emph{basho wa kujibiki ni yorimasu} — The places depend on drawing lots.
  \item \emph{Recall:} say the Japanese for fate, then the Japanese for a lottery, then say \emph{kujibiki} again
\end{itemize}

\begin{tcolorbox}[breakable,colback=teal!4,colframe=teal!35!black,title={Before you move on}]
What does \ja{くじびき} mean? (\textquotedblleft{}Drawing lots\textquotedblright{}.) Say it once more.
\end{tcolorbox}
\section[\texorpdfstring{\ja{ちゅうせんする} (chūsen suru)}{ちゅうせんする (chūsen suru)}]{\ja{ちゅうせんする} (chūsen suru) — to hold a draw; to draw lots}

\label{lesson:JA-C509-chuusen-suru}



\begin{quote}
Before the new one: say the Japanese for a lottery, then the Japanese for drawing lots.
\end{quote}

\subsection*{You'll want to know: \ja{ちゅうせんする}}
\textbf{\ja{ちゅうせんする}} — \emph{chūsen suru} — \textquotedblleft{}to hold a draw; to draw lots\textquotedblright{}.

\begin{quote}
\textbf{\ja{もし} \ja{せんにん} \ja{きたら、ちゅうせんします。}} — \emph{moshi sennin kitara, chūsenshimasu} — \textquotedblleft{}If a thousand people come, we will hold a draw.\textquotedblright{}
\end{quote}

Polite \textbf{\ja{ちゅうせんします}}, \emph{chūsenshimasu}; a \ja{する} verb. Past (polite): \textbf{\ja{みせが} \ja{ちゅうせんしました。}}, \emph{mise ga chūsenshimashita}.

\begin{cousinweb}
\textbf{\ja{ちゅうせんする}} (\emph{chūsen suru}) is the Sino-Japanese (kango) noun \emph{chūsen} with the native verb \emph{suru}, \textquotedblleft{}to do\textquotedblright{}: two Chinese characters read with their on'yomi, readings borrowed from Chinese. They are \emph{chū}, \textquotedblleft{}pluck\textquotedblright{} (Middle Chinese \emph{trhjuw}) and \emph{sen}, \textquotedblleft{}elect\textquotedblright{} (\emph{sjwenX}) (KANJIDIC2; Middle Chinese in Baxter's transcription).
\end{cousinweb}

\begin{grammarlens}[title={What you've built}]
One more word for saying what happens if.
\end{grammarlens}

\subsection*{Your turn}
\begin{itemize}
\raggedright
  \item \emph{Say it:} \emph{chūsen suru}
  \item \emph{Say it:} \emph{chūsen suru}, once more
  \item \emph{Say it:} \emph{moshi sennin kitara, chūsenshimasu} — If a thousand people come, we will hold a draw.
  \item \emph{Recall:} say the Japanese for a lottery, then the Japanese for drawing lots, then say \emph{chūsen suru} again
\end{itemize}

\begin{tcolorbox}[breakable,colback=teal!4,colframe=teal!35!black,title={Before you move on}]
What does \ja{ちゅうせんする} mean? (\textquotedblleft{}To hold a draw; to draw lots\textquotedblright{}.) Say it once more.
\end{tcolorbox}
\section[\texorpdfstring{\ja{かくりつ} (kakuritsu)}{かくりつ (kakuritsu)}]{\ja{かくりつ} (kakuritsu) — probability; the odds}

\label{lesson:JA-C509-kakuritsu}



\begin{quote}
Before the new one: say the Japanese for drawing lots, then the Japanese for to hold a draw.
\end{quote}

\subsection*{You'll want to know: \ja{かくりつ}}
\textbf{\ja{かくりつ}} — \emph{kakuritsu} — \textquotedblleft{}probability; the odds\textquotedblright{}.

\begin{quote}
\textbf{\ja{しあいは} \ja{あめの} \ja{かくりつに} \ja{よります。}} — \emph{shiai wa ame no kakuritsu ni yorimasu} — \textquotedblleft{}The match depends on the chance of rain.\textquotedblright{}
\end{quote}

\begin{cousinweb}
\textbf{\ja{かくりつ}} (\emph{kakuritsu}) is Sino-Japanese (kango): two Chinese characters read with their on'yomi, readings borrowed from Chinese. They are \emph{kaku}, \textquotedblleft{}assurance\textquotedblright{} (Middle Chinese \emph{khaewk}) and \emph{ritsu}, \textquotedblleft{}ratio\textquotedblright{} (KANJIDIC2; Middle Chinese in Baxter's transcription).
\end{cousinweb}

\begin{grammarlens}[title={What you've built}]
One more word for saying what happens if.
\end{grammarlens}

\subsection*{Your turn}
\begin{itemize}
\raggedright
  \item \emph{Say it:} \emph{kakuritsu}
  \item \emph{Say it:} \emph{kakuritsu}, once more
  \item \emph{Say it:} \emph{shiai wa ame no kakuritsu ni yorimasu} — The match depends on the chance of rain.
  \item \emph{Recall:} say the Japanese for drawing lots, then the Japanese for to hold a draw, then say \emph{kakuritsu} again
\end{itemize}

\begin{tcolorbox}[breakable,colback=teal!4,colframe=teal!35!black,title={Before you move on}]
What does \ja{かくりつ} mean? (\textquotedblleft{}Probability; the odds\textquotedblright{}.) Say it once more.
\end{tcolorbox}
\section[\texorpdfstring{\ja{あてる} (ateru)}{あてる (ateru)}]{\ja{あてる} (ateru) — to hit (the mark); to guess right}

\label{lesson:JA-C509-ateru}



\begin{quote}
Before the new one: say the Japanese for to hold a draw, then the Japanese for probability.
\end{quote}

\subsection*{You'll want to know: \ja{あてる}}
\textbf{\ja{あてる}} — \emph{ateru} — \textquotedblleft{}to hit (the mark); to guess right\textquotedblright{}.

\begin{quote}
\textbf{\ja{もし} \ja{こたえを} \ja{あてたら、おかしを} \ja{かいます。}} — \emph{moshi kotae o atetara, okashi o kaimasu} — \textquotedblleft{}If you guess the answer, I will buy you some sweets.\textquotedblright{}
\end{quote}

Polite \textbf{\ja{あてます}}, \emph{atemasu}; an ichidan verb. Past (polite): \textbf{\ja{おとうとが} \ja{こたえを} \ja{あてました。}}, \emph{otōto ga kotae o atemashita}.

\begin{cousinweb}
\textbf{\ja{あてる}} (\emph{ateru}) is a native Japanese word (wago), inherited rather than borrowed from Chinese; its kanji was attached later as a kun'yomi, and its older origin is not settled (JMdict, KANJIDIC2).
\end{cousinweb}

\begin{grammarlens}[title={What you've built}]
One more word for saying what happens if.
\end{grammarlens}

\subsection*{Your turn}
\begin{itemize}
\raggedright
  \item \emph{Say it:} \emph{ateru}
  \item \emph{Say it:} \emph{ateru}, once more
  \item \emph{Say it:} \emph{moshi kotae o atetara, okashi o kaimasu} — If you guess the answer, I will buy you some sweets.
  \item \emph{Recall:} say the Japanese for to hold a draw, then the Japanese for probability, then say \emph{ateru} again
\end{itemize}

\begin{tcolorbox}[breakable,colback=teal!4,colframe=teal!35!black,title={Before you move on}]
What does \ja{あてる} mean? (\textquotedblleft{}To hit (the mark); to guess right\textquotedblright{}.) Say it once more.
\end{tcolorbox}
\section[(dialogue)]{Looking back: from Books and Art to Lotteries and Odds}

\label{lesson:JA-R509-review-504-509}



\begin{quote}
Say the Japanese for literature and a poet.
\end{quote}

\subsection*{Your turn}
Say these aloud:

\begin{itemize}
\raggedright
  \item literature, a poet, an art gallery, an artist, a picture book
  \item a sumo wrestler, a sumo ring, a sumo referee, a bow (for arrows), to dodge
  \item a sport, a championship, a victory, a defeat, to defeat
  \item a hope, to hope for, to aspire to, an ideal, to want (said of someone else)
  \item luck, good fortune, a misfortune, luckily, fate
  \item a lottery, drawing lots, to hold a draw, probability, to hit (the mark)
\end{itemize}

\begin{tcolorbox}[breakable,colback=teal!4,colframe=teal!35!black,title={Before you move on}]
Literature? (\textbf{\ja{ぶんがく}}, \emph{bungaku}.) A poet? (\textbf{\ja{しじん}}, \emph{shijin}.) An art gallery? (\textbf{\ja{がろう}}, \emph{garō}.) An artist? (\textbf{\ja{げいじゅつか}}, \emph{geijutsuka}.) A picture book? (\textbf{\ja{えほん}}, \emph{ehon}.) A sumo wrestler? (\textbf{\ja{りきし}}, \emph{rikishi}.) A sumo ring? (\textbf{\ja{どひょう}}, \emph{dohyō}.) A sumo referee? (\textbf{\ja{ぎょうじ}}, \emph{gyōji}.) A bow (for arrows)? (\textbf{\ja{ゆみ}}, \emph{yumi}.) To dodge? (\textbf{\ja{よける}}, \emph{yokeru}.) A sport? (\textbf{\ja{きょうぎ}}, \emph{kyōgi}.) A championship? (\textbf{\ja{せんしゅけん}}, \emph{senshuken}.) A victory? (\textbf{\ja{しょうり}}, \emph{shōri}.) A defeat? (\textbf{\ja{はいぼく}}, \emph{haiboku}.) To defeat? (\textbf{\ja{まかす}}, \emph{makasu}.) A hope? (\textbf{\ja{きぼう}}, \emph{kibō}.) To hope for? (\textbf{\ja{のぞむ}}, \emph{nozomu}.) To aspire to? (\textbf{\ja{しぼうする}}, \emph{shibō suru}.) An ideal? (\textbf{\ja{りそう}}, \emph{risō}.) To want (said of someone else)? (\textbf{\ja{ほしがる}}, \emph{hoshigaru}.) Luck? (\textbf{\ja{うん}}, \emph{un}.) Good fortune? (\textbf{\ja{こううん}}, \emph{kōun}.) A misfortune? (\textbf{\ja{さいなん}}, \emph{sainan}.) Luckily? (\textbf{\ja{うんよく}}, \emph{un'yoku}.) Fate? (\textbf{\ja{うんめい}}, \emph{unmei}.) A lottery? (\textbf{\ja{たからくじ}}, \emph{takarakuji}.) Drawing lots? (\textbf{\ja{くじびき}}, \emph{kujibiki}.) To hold a draw? (\textbf{\ja{ちゅうせんする}}, \emph{chūsen suru}.) Probability? (\textbf{\ja{かくりつ}}, \emph{kakuritsu}.) To hit (the mark)? (\textbf{\ja{あてる}}, \emph{ateru}.)
\end{tcolorbox}
